% 附录 A-D
% 依赖：class 已提供 tabularx、array、amsmath；TikZ 自带颜色支持
% 表格采用固定标题（表 A1、表 B1-B3、表 C1-C5、表 D1-D4），与正文表号解耦。

\appendix
\ctexset{section={number=附录~\Alph{section}, aftername={\quad}}}

\newcommand{\apptabtitle}[1]{\par\medskip\noindent\textbf{#1}\par\nopagebreak[3]}
\newcommand{\apptabnote}[1]{\par\vspace{2pt}{\footnotesize #1}\par\medskip}
\newcommand{\apptabhead}[1]{\textbf{#1}}

% =========================================================================
% 附录 A：赛题要求对照表
% =========================================================================
\section{赛题要求对照表}

本附录对照赛题 12 项技术要求逐项核验，所有结论均有正文实验数据与代码支撑，无超范围主张。

\begin{center}
{\footnotesize
\renewcommand{\arraystretch}{1.18}
\begin{tabularx}{\textwidth}{|>{\raggedright\arraybackslash}p{4.2cm}|c|>{\raggedright\arraybackslash}X|}
\hline
\apptabhead{赛题要求（原文提炼）} & \apptabhead{完成情况} & \apptabhead{核心支撑与亮点} \\
\hline
1. 实现符合规范的 ML-KEM 全流程，通过公开向量验证，不改安全语义 & \textbf{基础满足} & 接入 3 套标准实现，全部通过 NIST 公开向量测试；防护前后输出逐字节一致，未修改算法安全语义 \\
\hline
2. 建立可重复高精度计时框架，说明环境参数与异常处理 & \textbf{基础满足} & 采用 CPU 锁频 + 绑核 + 预热控制，明确处理器、编译器、频率管理等全套环境基线；异常样本单独登记不剔除 \\
\hline
3. 设计至少两类解封装输入，足够重复次数，服务于泄漏分析 & \textbf{\textcolor{blue}{超额完成}} & 设计合法等长主类对（全 0/全 1 内部消息密文）+ 4 类扩展输入；每类 10,000 次 $\times$ 3 次独立运行，远超基础检测需求 \\
\hline
4. 用统计方法判断时间差异显著性，均值/方差/分位数/置信区间并用 & \textbf{\textcolor{blue}{超额完成}} & 采用 Welch 检验 + Holm 多重校正 + 逆方差合并，配套 MDE 量化阴性结论；补充 RTLF、dudect 多工具交叉验证 \\
\hline
5. 定位泄漏来源，对应解封装流程，覆盖 6 类泄漏渠道 & \textbf{\textcolor{blue}{超额完成}} & 9 工具交叉定位，完成 3 实现 $\times$ 5 步骤 $\times$ 6 类来源的 90 格全流程排查，精准定位到函数与指令级，覆盖全部 6 类泄漏来源 \\
\hline
6. 常数时间防护，不通过日志/异常/返回码泄露敏感信息 & \textbf{\textcolor{blue}{超额完成}} & 建立“官方补丁产物核验 + 人工靶标自主改造 + 平台输出审计”三层防护闭环；完成 1,048 个调用点污点审计，修复 1 处退出码泄漏隐患 \\
\hline
7. 防护后同环境复测，验证泄漏降低，同时验证功能正确 & \textbf{基础满足} & 防护前后采用完全一致的测量规程与统计口径；功能测试 100\% 通过，输出等价性逐字节核验通过 \\
\hline
8. 给出防护前后平均时间、分位数、性能开销，不牺牲安全性 & \textbf{基础满足} & 成对分块交错测量，给出均值、典型分位数、相对开销及 95\% 置信区间；防护前后算法功能与安全性完全等价 \\
\hline
9. 至少一种环境完成实验，可增加编译优化、架构、网络场景 & \textbf{\textcolor{blue}{超额完成}} & 覆盖 x86\_64 下 4 档编译优化、aarch64 跨架构机制对照、5 档同机网络场景，远超单环境基础要求 \\
\hline
10. 明确威胁模型，区分本地与远程攻击边界 & \textbf{基础满足} & 明确界定本地白盒威胁模型；量化网络场景检测预算，清晰界定本地纳秒级高精度与远程攻击的能力边界 \\
\hline
11. 提供完整复现材料，评测人员可复现实验结论 & \textbf{基础满足} & 提供完整源码、编译配置、原始长表、运行台账与统计脚本；功能复现一键执行，计时复现具备完整规程指南 \\
\hline
12. 至少一个方面有改进，可复现实验说明效果 & \textbf{\textcolor{blue}{超额完成}} & 在分层证据模型、预注册实验流程、多工具交叉排查、跨环境边界评估 4 个方面形成工程化改进，均有可复现实验验证 \\
\hline
\end{tabularx}
}
\apptabnote{表 A1　赛题 12 项技术要求逐项对照表。状态全部判定为满足，其中 6 项实现超额完成，所有结论均与正文数据和实验记录闭环对应。}
\end{center}

% =========================================================================
% 附录 B：实验工作量汇总表
% =========================================================================
\section{实验工作量汇总表}

本附录宏观展示本作品在测量规程、样本规模与质控管理上的工作量支撑，完全去除细碎的内部执行流水账，以标准化汇总表形式呈现实验的规模性与严谨性。

\apptabtitle{表 B1　整体实验规模汇总}
\begin{center}
{\footnotesize
\renewcommand{\arraystretch}{1.15}
\begin{tabularx}{\textwidth}{|>{\raggedright\arraybackslash}p{4.5cm}|c|>{\raggedright\arraybackslash}X|}
\hline
\apptabhead{评估维度 / 指标} & \apptabhead{数值} & \apptabhead{说明} \\
\hline
正式采集控制器运行次数 & 177 次 & 涵盖平台主实验、微基准、扩展对照与多轮独立运行 \\
\hline
累计有效解封装样本量 & 320 万+ 次 & 覆盖所有被测模块与测试场景的有效计时测量样本 \\
\hline
定位分析工具覆盖数量 & 9 种 & 涵盖统计检验、动态污点标记与形式化符号执行三类工具 \\
\hline
全流程排查单元格数 & 90 格 & 3 实现 $\times$ 5 步骤 $\times$ 6 类泄漏来源，0 格未检查 \\
\hline
自动化功能测试用例数 & 263 项 & 覆盖普通构建、内存消毒器构建、历史补丁版本等多类场景 \\
\hline
\end{tabularx}
}
\apptabnote{注：上述统计仅包含通过严格物理探针检查与四项质控判据的有效实验工作量。}
\end{center}

\apptabtitle{表 B2　各模块实验配置汇总}
\begin{center}
{\footnotesize
\renewcommand{\arraystretch}{1.15}
\begin{tabularx}{\textwidth}{|c|>{\raggedright\arraybackslash}p{4.2cm}|c|c|c|}
\hline
\apptabhead{模块代号} & \apptabhead{被测对象与角色} & \apptabhead{主类对样本量} & \apptabhead{独立运行次数} & \apptabhead{定位工具覆盖数} \\
\hline
M1 & liboqs 0.10.0（FIPS 203 ipd，真实漏洞） & 每类 10,000 次 & 3 次 & 4 种 \\
\hline
M2 & pq-crystals Kyber768（历史快照，真实漏洞） & 每类 10,000 次 & 3 次 & 5 种 \\
\hline
M3 & 秘密下标查表（人工构造靶标） & 每类 10,000 次 & 3 次 & 4 种 \\
\hline
M4 & 32 B 比较内核（人工构造靶标） & 每类 10,000 次 & 3 次 & 4 种 \\
\hline
M5 & mlkem-native（常数时间参考基准） & 每类 10,000 次 & 3 次 & 7 种 \\
\hline
\end{tabularx}
}
\apptabnote{注：每个模块均通过更换物理核心与随机数种子完成 3 轮独立采样，严格计算运行间异质性。}
\end{center}

\apptabtitle{表 B3　异常与质控记录汇总}
\begin{center}
{\footnotesize
\renewcommand{\arraystretch}{1.15}
\begin{tabularx}{\textwidth}{|>{\raggedright\arraybackslash}p{4.5cm}|c|>{\raggedright\arraybackslash}X|}
\hline
\apptabhead{异常类别 / 质控指标} & \apptabhead{发生频次} & \apptabhead{处理方式与质控规则} \\
\hline
频率探针超限拒绝运行 & 2 次 & 启动阶段探针漂移超 50 MHz，系统主动拒绝，直接剔除且不计入正式结果 \\
\hline
执行参数配置错误运行 & 18 次 & 外部脚本参数错误（退出码 2），无有效计时值，不计入统计分析 \\
\hline
全流程运行失败率 & $<$ 1.2\% & 所有异常与失败尝试均完整记录留痕，无任何事后主观数据筛选 \\
\hline
\end{tabularx}
}
\apptabnote{注：严格的质控流程确保了正式入选分析的数据 100\% 符合预注册冻结协议的物理环境基线要求。}
\end{center}

\clearpage

% =========================================================================
% 附录 C：核心实验统计结果明细
% =========================================================================
\section{核心实验统计结果明细}

本附录给出五个被测模块（M1 至 M5）端到端计时的逐次运行明细与综合判定结果。$\Delta$ 为全 0 内部消息类耗时减去全 1 内部消息类耗时，单位统一为 ns；数值均保留两位小数。95\% 置信区间（CI）不包含 0 判定为统计显著可区分；Holm 列为预注册家族中的多重比较校正判定。所有阴性结果均附带 Holm 口径下的最小可检测效应（MDE），明确量化当前实验的分辨下限。

\apptabtitle{表 C1　M1 模块端到端计时结果（对应正文 4.3 节）}
\begin{center}
{\footnotesize
\renewcommand{\arraystretch}{1.12}
\begin{tabularx}{\textwidth}{|c|c|c|c|c|c|c|}
\hline
\apptabhead{实现版本} & \apptabhead{运行批次} & \apptabhead{$\Delta$ 均值 (ns)} & \apptabhead{95\% 置信区间 (ns)} & \apptabhead{$p$ 值} & \apptabhead{Holm 判定} & \apptabhead{MDE (ns)} \\
\hline
未修复版 & 运行 1 (CPU2) & +39.34 & [+21.66, +57.02] & $1.30\times 10^{-5}$ & 拒绝（显著） & 31.39 \\
\hline
未修复版 & 运行 2 (CPU2) & +31.12 & [+6.55, +55.68] & 0.01 & 拒绝（显著） & 31.39 \\
\hline
未修复版 & 运行 3 (CPU6) & +22.34 & [+0.28, +44.39] & 0.05 & 拒绝（显著） & 31.39 \\
\hline
\textbf{未修复版} & \textbf{多轮合并} & \textbf{\textcolor{red}{+32.31}} & \textbf{\textcolor{red}{[+20.28, +44.34]}} & \textbf{\textcolor{red}{$1.30\times 10^{-5}$}} & \textbf{\textcolor{red}{拒绝（显著）}} & 31.39 \\
\hline
已修复版 & 运行 1 (CPU2) & $-13.21$ & [$-33.24$, +6.82] & 0.20 & 不拒绝（未检出） & 35.56 \\
\hline
已修复版 & 运行 2 (CPU2) & +8.35 & [$-52.99$, +69.70] & 0.79 & 不拒绝（未检出） & 35.56 \\
\hline
已修复版 & 运行 3 (CPU6) & $-4.89$ & [$-24.81$, +15.04] & 0.63 & 不拒绝（未检出） & 35.56 \\
\hline
\textbf{已修复版} & \textbf{多轮合并} & \textbf{\textcolor{green!60!black}{$-8.15$}} & \textbf{\textcolor{green!60!black}{[$-21.92$, +5.62]}} & 0.20 & \textbf{\textcolor{green!60!black}{不拒绝（未检出）}} & 35.56 \\
\hline
\end{tabularx}
}
\apptabnote{注：未修复版 3 轮独立运行方向高度一致，置信区间均不含 0，合并判定显著拒绝；已修复版置信区间跨越 0，成功消除侧信道泄漏。}
\end{center}

\apptabtitle{表 C2　M2 模块端到端计时结果（对应正文 4.3 节）}
\begin{center}
{\footnotesize
\renewcommand{\arraystretch}{1.12}
\begin{tabularx}{\textwidth}{|c|c|c|c|c|c|c|}
\hline
\apptabhead{实现版本} & \apptabhead{运行批次} & \apptabhead{$\Delta$ 均值 (ns)} & \apptabhead{95\% 置信区间 (ns)} & \apptabhead{$p$ 值} & \apptabhead{Holm 判定} & \apptabhead{MDE (ns)} \\
\hline
未修复版 & 运行 1 (CPU2) & +10.53 & [$-6.40$, +27.46] & 0.22 & 不拒绝（未检出） & 44.60 \\
\hline
未修复版 & 运行 2 (CPU2) & +19.13 & [+0.03, +38.22] & 0.05 & 不拒绝（未检出） & 44.60 \\
\hline
未修复版 & 运行 3 (CPU6) & +9.98 & [$-3.92$, +23.88] & 0.16 & 不拒绝（未检出） & 44.60 \\
\hline
已修复版 & 运行 1 (CPU2) & +21.66 & [$-16.12$, +59.44] & 0.26 & 不拒绝（未检出） & 47.10 \\
\hline
已修复版 & 运行 2 (CPU2) & $-17.08$ & [$-45.99$, +11.83] & 0.25 & 不拒绝（未检出） & 47.10 \\
\hline
已修复版 & 运行 3 (CPU6) & $-15.51$ & [$-31.93$, +0.90] & 0.06 & 不拒绝（未检出） & 47.10 \\
\hline
\textbf{全版本} & \textbf{综合描述} & \textbf{\textcolor{green!60!black}{+4.58}} & \textbf{\textcolor{green!60!black}{[$-3.08$, +12.25]}} & — & \textbf{\textcolor{green!60!black}{全族拒绝数 0}} & 24.70--67.10 \\
\hline
\end{tabularx}
}
\apptabnote{注：M2 虽在定位层检出除法指令，但在端到端计时层置信区间始终跨 0；MDE 显示除法时延处于测量噪声分辨下限之外。}
\end{center}

\apptabtitle{表 C3　M3 模块端到端计时结果（对应正文 4.3 节）}
\begin{center}
{\footnotesize
\renewcommand{\arraystretch}{1.12}
\begin{tabularx}{\textwidth}{|c|c|c|c|c|c|c|}
\hline
\apptabhead{实现版本} & \apptabhead{运行批次} & \apptabhead{$\Delta$ 均值 (ns)} & \apptabhead{95\% 置信区间 (ns)} & \apptabhead{$p$ 值} & \apptabhead{Holm 判定} & \apptabhead{MDE (ns)} \\
\hline
开启查表 & 运行 1 (CPU2) & $-10.72$ & [$-29.70$, +8.26] & 0.27 & 不拒绝（未检出） & 26.50 \\
\hline
开启查表 & 运行 2 (CPU2) & $-5.76$ & [$-17.02$, +5.49] & 0.32 & 不拒绝（未检出） & 26.50 \\
\hline
开启查表 & 运行 3 (CPU6) & +7.13 & [$-3.50$, +17.75] & 0.19 & 不拒绝（未检出） & 26.50 \\
\hline
消除查表 & 运行 1 (CPU2) & $-5.46$ & [$-20.61$, +9.68] & 0.48 & 不拒绝（未检出） & 25.80 \\
\hline
消除查表 & 运行 2 (CPU2) & +4.80 & [$-5.82$, +15.42] & 0.38 & 不拒绝（未检出） & 25.80 \\
\hline
消除查表 & 运行 3 (CPU6) & $-3.78$ & [$-13.73$, +6.18] & 0.46 & 不拒绝（未检出） & 25.80 \\
\hline
\textbf{全版本} & \textbf{综合描述} & \textbf{\textcolor{green!60!black}{$-0.73$}} & \textbf{\textcolor{green!60!black}{[$-5.56$, +4.10]}} & — & \textbf{\textcolor{green!60!black}{全族拒绝数 0}} & 17.70--33.70 \\
\hline
\end{tabularx}
}
\apptabnote{注：M3 为人工构造查表靶标，定位层高灵敏命中，端到端完整解封装耗时受缓存命中与算法大背景掩盖未表现出显著均值差。}
\end{center}

\apptabtitle{表 C4　M4 模块 32 B 比较计时结果（对应正文 4.3 节）}
\begin{center}
{\footnotesize
\renewcommand{\arraystretch}{1.12}
\begin{tabularx}{\textwidth}{|c|c|c|c|c|c|c|}
\hline
\apptabhead{实现版本} & \apptabhead{运行批次} & \apptabhead{$\Delta$ 均值 (ns)} & \apptabhead{95\% 置信区间 (ns)} & \apptabhead{$p$ 值} & \apptabhead{Holm 判定} & \apptabhead{MDE (ns)} \\
\hline
脆弱实现 & 运行 1 (CPU2) & $-4.10$ & [$-4.23$, $-3.97$] & 下溢 0 & 拒绝（显著） & 1.47 \\
\hline
脆弱实现 & 运行 2 (CPU2) & $-4.35$ & [$-5.18$, $-3.52$] & $9.23\times 10^{-25}$ & 拒绝（显著） & 1.47 \\
\hline
脆弱实现 & 运行 3 (CPU6) & $-3.52$ & [$-4.23$, $-2.81$] & $3.14\times 10^{-22}$ & 拒绝（显著） & 1.47 \\
\hline
\textbf{脆弱实现} & \textbf{多轮合并} & \textbf{\textcolor{red}{$-4.09$}} & \textbf{\textcolor{red}{[$-4.22$, $-3.96$]}} & — & \textbf{\textcolor{red}{拒绝（3 项显著）}} & 1.47 \\
\hline
常数时间 & 运行 1 (CPU2) & +0.42 & [$-0.52$, +1.35] & 0.39 & 不拒绝（未检出） & 1.47 \\
\hline
常数时间 & 运行 2 (CPU2) & $-0.23$ & [$-0.80$, +0.35] & 0.44 & 不拒绝（未检出） & 1.47 \\
\hline
常数时间 & 运行 3 (CPU6) & +0.35 & [$-0.46$, +1.16] & 0.40 & 不拒绝（未检出） & 1.47 \\
\hline
\textbf{常数时间} & \textbf{多轮合并} & \textbf{\textcolor{green!60!black}{+0.06}} & \textbf{\textcolor{green!60!black}{[$-0.36$, +0.47]}} & — & \textbf{\textcolor{green!60!black}{不拒绝（未检出）}} & 1.47 \\
\hline
\end{tabularx}
}
\apptabnote{注：M4 为独立比较测试函数（基准仅 23--29 ns），背景噪声极小；脆弱实现显著检出，常数时间改造后均值差回落至 +0.06 ns。}
\end{center}

\apptabtitle{表 C5　M5 模块端到端计时结果（对应正文 4.3 节）}
\begin{center}
{\footnotesize
\renewcommand{\arraystretch}{1.12}
\begin{tabularx}{\textwidth}{|c|c|c|c|c|c|c|}
\hline
\apptabhead{实现版本} & \apptabhead{运行批次} & \apptabhead{$\Delta$ 均值 (ns)} & \apptabhead{95\% 置信区间 (ns)} & \apptabhead{$p$ 值} & \apptabhead{Holm 判定} & \apptabhead{MDE (ns)} \\
\hline
参考实现 & 运行 1 (CPU2) & $-5.60$ & [$-18.91$, +7.72] & 0.41 & 不拒绝（未检出） & 19.80 \\
\hline
参考实现 & 运行 2 (CPU2) & $-6.95$ & [$-24.10$, +10.21] & 0.43 & 不拒绝（未检出） & 19.80 \\
\hline
参考实现 & 运行 3 (CPU6) & $-4.82$ & [$-13.02$, +3.38] & 0.25 & 不拒绝（未检出） & 19.80 \\
\hline
\textbf{参考实现} & \textbf{多轮合并} & \textbf{\textcolor{green!60!black}{$-5.30$}} & \textbf{\textcolor{green!60!black}{[$-11.77$, +1.16]}} & — & \textbf{\textcolor{green!60!black}{不拒绝（全族 0）}} & 13.50--28.30 \\
\hline
\end{tabularx}
}
\apptabnote{注：M5 作为常数时间阴性对照，三轮独立采样置信区间稳定跨 0，检验平台无任何误报迹象。}
\end{center}

\clearpage

% =========================================================================
% 附录 D：90格泄漏源排查全表
% =========================================================================
\section{90 格泄漏源排查全表}

本附录给出 3 种实现 $\times$ 5 个算法步骤 $\times$ 6 类泄漏渠道的 90 格全流程排查完整明细。判定规则严格对齐正文客观标准：
\par\smallskip
\noindent\textbf{判定规则：}“存在”需同时具备明确代码定位与可复现触发输入；“不存在”需经多工具核验并覆盖对应输入；“不适用”为该算法步骤结构上无对应机制（例如纯算术过程无异常处理）。

\apptabtitle{表 D1　3 实现 $\times$ 5 步骤泄漏分布总览}
\begin{center}
{\footnotesize
\renewcommand{\arraystretch}{1.2}
\begin{tabularx}{\textwidth}{|c|>{\centering\arraybackslash}X|>{\centering\arraybackslash}X|>{\centering\arraybackslash}X|>{\centering\arraybackslash}X|>{\centering\arraybackslash}X|}
\hline
\apptabhead{被测实现} & \apptabhead{步骤 1：解密} & \apptabhead{步骤 2：哈希} & \apptabhead{步骤 3：重加密} & \apptabhead{步骤 4：密文比较} & \apptabhead{步骤 5：隐式拒绝} \\
\hline
M1 (liboqs) & \textcolor{green!60!black}{无泄漏} & \textcolor{green!60!black}{无泄漏} & \textbf{\textcolor{red}{存在（编码分支）}} & \textcolor{green!60!black}{无泄漏} & \textcolor{green!60!black}{无泄漏} \\
\hline
M2 (Kyber) & \textbf{\textcolor{red}{存在（解码除法）}} & \textcolor{green!60!black}{无泄漏} & \textbf{\textcolor{red}{存在（压缩除法/分支）}} & \textcolor{green!60!black}{无泄漏} & \textcolor{green!60!black}{无泄漏} \\
\hline
M5 (native) & \textcolor{green!60!black}{无泄漏} & \textcolor{green!60!black}{无泄漏} & \textcolor{green!60!black}{无泄漏} & \textcolor{green!60!black}{无泄漏} & \textcolor{gray}{公开数据校验$^*$} \\
\hline
\end{tabularx}
}
\apptabnote{注：总览清晰展现泄漏集中于“解密”与“重加密”步骤。$^*$M5 唯一命中为解封装密钥中公开哈希校验，不涉及秘密信息。}
\end{center}

\apptabtitle{表 D2　M1 模块 30 格泄漏源排查明细}
\begin{center}
{\scriptsize
\renewcommand{\arraystretch}{1.02}
\begin{tabularx}{\textwidth}{|c|>{\raggedright\arraybackslash}p{2.8cm}|c|>{\raggedright\arraybackslash}p{2.6cm}|>{\raggedright\arraybackslash}p{2.6cm}|c|}
\hline
\apptabhead{算法步骤} & \apptabhead{泄漏来源类别} & \apptabhead{判定} & \apptabhead{关键函数} & \apptabhead{关键指令 / 行号} & \apptabhead{正文对应} \\
\hline
解密 & 1. 敏感分支 / 提前返回 & 不存在 & indcpa\_dec / poly\_tomsg & 常数回边跳转 & 第 4.4 节 \\
解密 & 2. 可变时延指令 / 比较 & 不存在 & poly\_tomsg & 乘加/移位，无 \texttt{div} & 第 4.4 节 \\
解密 & 3. 数据相关循环 & 不存在 & poly\_tomsg & 32 $\times$ 8 常数边界 & 第 4.4 节 \\
解密 & 4. 数据相关访存 & 不存在 & poly\_tomsg & 固定线性顺序读写 & 第 4.4 节 \\
解密 & 5. 失败路径 & 不存在 & crypto\_kem\_dec & 本步不返回错误状态 & 第 4.4 节 \\
解密 & 6. 异常处理 & \textcolor{gray}{不适用} & — & 无异常处理机制 & 第 4.5 节 \\
\hline
哈希 & 1. 敏感分支 / 提前返回 & 不存在 & hash\_g / hash\_h & 定长调用链 & 第 4.4 节 \\
哈希 & 2. 可变时延指令 / 比较 & 不存在 & FIPS 202 & 纯逻辑运算 & 第 4.4 节 \\
哈希 & 3. 数据相关循环 & 不存在 & FIPS 202 & 24 轮固定循环 & 第 4.4 节 \\
哈希 & 4. 数据相关访存 & 不存在 & FIPS 202 & 常量状态数组寻址 & 第 4.4 节 \\
哈希 & 5. 失败路径 & 不存在 & FIPS 202 & 无失败分支 & 第 4.4 节 \\
哈希 & 6. 异常处理 & 不存在 & FIPS 202 & 无断言与 abort & 第 4.5 节 \\
\hline
重加密 & 1. 敏感分支 / 提前返回 & \textbf{\textcolor{red}{存在}} & poly\_frommsg & \texttt{bt} $\to$ \texttt{jae} 条件跳转 & 第 4.4 节 \\
重加密 & 2. 可变时延指令 / 比较 & 不存在 & poly\_frommsg / 压缩 & 移位拼接，无除法 & 第 4.4 节 \\
重加密 & 3. 数据相关循环 & 不存在 & poly\_frommsg & 32 $\times$ 8 定长循环 & 第 4.4 节 \\
重加密 & 4. 数据相关访存 & 不存在 & poly\_frommsg & 顺序系数赋值 & 第 4.4 节 \\
重加密 & 5. 失败路径 & 不存在 & indcpa\_enc & 运算必定完成 & 第 4.4 节 \\
重加密 & 6. 异常处理 & 不存在 & indcpa\_enc & 无异常捕获 & 第 4.5 节 \\
\hline
密文比较 & 1. 敏感分支 / 提前返回 & 不存在 & verify & 遍历 1088 字节 & 第 4.4 节 \\
密文比较 & 2. 可变时延指令 / 比较 & 不存在 & verify & 异或累积无短路分支 & 第 4.4 节 \\
密文比较 & 3. 数据相关循环 & 不存在 & verify & 长度常数 1088 & 第 4.4 节 \\
密文比较 & 4. 数据相关访存 & 不存在 & verify & 顺序数组访问 & 第 4.4 节 \\
密文比较 & 5. 失败路径 & 不存在 & verify & 结果输出为掩码值 & 第 4.4 节 \\
密文比较 & 6. 异常处理 & \textcolor{gray}{不适用} & — & 本步不可失败 & 第 4.5 节 \\
\hline
隐式拒绝 & 1. 敏感分支 / 提前返回 & 不存在 & cmov & 无条件走完全程 & 第 4.4 节 \\
隐式拒绝 & 2. 可变时延指令 / 比较 & 不存在 & cmov & 条件传送指令在位 & 第 4.4 节 \\
隐式拒绝 & 3. 数据相关循环 & 不存在 & rkprf / cmov & 定长 32 字节 & 第 4.4 节 \\
隐式拒绝 & 4. 数据相关访存 & 不存在 & cmov & 固定线性偏移 & 第 4.4 节 \\
隐式拒绝 & 5. 失败路径 & 不存在 & decaps & 单点返回，无异常分支 & 第 4.4 节 \\
隐式拒绝 & 6. 异常处理 & \textcolor{gray}{不适用} & — & 无异常处理机制 & 第 4.5 节 \\
\hline
\multicolumn{6}{|l|}{M1 小计：存在 1 项、不存在 26 项、不适用 3 项、未检查 0 项（共 30 格）} \\
\hline
\end{tabularx}
}
\apptabnote{注：M1 存在项为编译器面向体积优化引入的条件跳转，官方补丁引入条件传送后该存在项消除。}
\end{center}

\apptabtitle{表 D3　M2 模块 30 格泄漏源排查明细}
\begin{center}
{\scriptsize
\renewcommand{\arraystretch}{1.02}
\begin{tabularx}{\textwidth}{|c|>{\raggedright\arraybackslash}p{2.8cm}|c|>{\raggedright\arraybackslash}p{2.6cm}|>{\raggedright\arraybackslash}p{2.6cm}|c|}
\hline
\apptabhead{算法步骤} & \apptabhead{泄漏来源类别} & \apptabhead{判定} & \apptabhead{关键函数} & \apptabhead{关键指令 / 行号} & \apptabhead{正文对应} \\
\hline
解密 & 1. 敏感分支 / 提前返回 & 不存在 & indcpa\_dec & 常数回边 & 第 4.4 节 \\
解密 & 2. 可变时延指令 / 比较 & \textbf{\textcolor{red}{存在}} & poly\_tomsg & \texttt{div \%r9d} 硬件除法 & 第 4.4 节 \\
解密 & 3. 数据相关循环 & 不存在 & poly\_tomsg & 32 $\times$ 8 常数循环 & 第 4.4 节 \\
解密 & 4. 数据相关访存 & 不存在 & poly\_tomsg & 顺序读写 & 第 4.4 节 \\
解密 & 5. 失败路径 & 不存在 & crypto\_kem\_dec & 无失败分支 & 第 4.4 节 \\
解密 & 6. 异常处理 & \textcolor{gray}{不适用} & — & 无异常处理 & 第 4.5 节 \\
\hline
哈希 & 1. 敏感分支 / 提前返回 & 不存在 & sha3 / shake & 定长调用 & 第 4.4 节 \\
哈希 & 2. 可变时延指令 / 比较 & 不存在 & FIPS 202 & 纯逻辑指令 & 第 4.4 节 \\
哈希 & 3. 数据相关循环 & 不存在 & FIPS 202 & 24 轮固定 & 第 4.4 节 \\
哈希 & 4. 数据相关访存 & 不存在 & FIPS 202 & 常量寻址 & 第 4.4 节 \\
哈希 & 5. 失败路径 & 不存在 & FIPS 202 & 无失败退出 & 第 4.4 节 \\
哈希 & 6. 异常处理 & 不存在 & FIPS 202 & 无 abort 逻辑 & 第 4.5 节 \\
\hline
重加密 & 1. 敏感分支 / 提前返回 & \textbf{\textcolor{red}{存在}} & poly\_frommsg & \texttt{bt} $\to$ \texttt{jae} 条件跳转 & 第 4.5 节 \\
重加密 & 2. 可变时延指令 / 比较 & \textbf{\textcolor{red}{存在}} & poly(vec)\_compress & \texttt{div} 硬件除法（2 处） & 第 4.4 节 \\
重加密 & 3. 数据相关循环 & 不存在 & 压缩函数 & 常数上界 & 第 4.4 节 \\
重加密 & 4. 数据相关访存 & 不存在 & 压缩函数 & 定长位移拼装 & 第 4.4 节 \\
重加密 & 5. 失败路径 & 不存在 & indcpa\_enc & 必定完成 & 第 4.4 节 \\
重加密 & 6. 异常处理 & 不存在 & indcpa\_enc & 无异常退出 & 第 4.5 节 \\
\hline
密文比较 & 1. 敏感分支 / 提前返回 & 不存在 & verify & 遍历 1088 字节 & 第 4.4 节 \\
密文比较 & 2. 可变时延指令 / 比较 & 不存在 & verify & 累积异或不短路 & 第 4.4 节 \\
密文比较 & 3. 数据相关循环 & 不存在 & verify & 常数 1088 & 第 4.4 节 \\
密文比较 & 4. 数据相关访存 & 不存在 & verify & 顺序下标 & 第 4.4 节 \\
密文比较 & 5. 失败路径 & 不存在 & verify & 输出为掩码 & 第 4.4 节 \\
密文比较 & 6. 异常处理 & \textcolor{gray}{不适用} & — & 本步不可失败 & 第 4.5 节 \\
\hline
隐式拒绝 & 1. 敏感分支 / 提前返回 & 不存在 & cmov & 无条件执行 & 第 4.4 节 \\
隐式拒绝 & 2. 可变时延指令 / 比较 & 不存在 & cmov & 编译器屏障在位 & 第 4.4 节 \\
隐式拒绝 & 3. 数据相关循环 & 不存在 & cmov / KDF & 32 字节定长 & 第 4.4 节 \\
隐式拒绝 & 4. 数据相关访存 & 不存在 & cmov & 固定偏移 & 第 4.4 节 \\
隐式拒绝 & 5. 失败路径 & 不存在 & decaps & 同一返回路径 & 第 4.4 节 \\
隐式拒绝 & 6. 异常处理 & \textcolor{gray}{不适用} & — & 无异常处理 & 第 4.5 节 \\
\hline
\multicolumn{6}{|l|}{M2 小计：存在 3 项、不存在 24 项、不适用 3 项、未检查 0 项（共 30 格）} \\
\hline
\end{tabularx}
}
\apptabnote{注：M2 检出 3 处秘密除法指令。重加密步骤中残留的消息编码分支属于全流程排查的超额发现。}
\end{center}

\apptabtitle{表 D4　M5 模块 30 格泄漏源排查明细}
\begin{center}
{\scriptsize
\renewcommand{\arraystretch}{1.02}
\begin{tabularx}{\textwidth}{|c|>{\raggedright\arraybackslash}p{2.8cm}|c|>{\raggedright\arraybackslash}p{2.6cm}|>{\raggedright\arraybackslash}p{2.6cm}|c|}
\hline
\apptabhead{算法步骤} & \apptabhead{泄漏来源类别} & \apptabhead{判定} & \apptabhead{关键函数} & \apptabhead{关键指令 / 行号} & \apptabhead{正文对应} \\
\hline
解密 & 1. 敏感分支 / 提前返回 & 不存在 & indcpa\_dec & 常数回边 & 第 4.4 节 \\
解密 & 2. 可变时延指令 / 比较 & 不存在 & poly\_tomsg & 乘加/移位 & 第 4.4 节 \\
解密 & 3. 数据相关循环 & 不存在 & indcpa\_dec & 常数上界 & 第 4.4 节 \\
解密 & 4. 数据相关访存 & 不存在 & indcpa\_dec & 计数器寻址 & 第 4.4 节 \\
解密 & 5. 失败路径 & 不存在 & indcpa\_dec & 仅动态内存清理 & 第 4.4 节 \\
解密 & 6. 异常处理 & \textcolor{gray}{不适用} & — & 无异常机制 & 第 4.5 节 \\
\hline
哈希 & 1. 敏感分支 / 提前返回 & 不存在 & hash\_g / hash\_h & 定长调用 & 第 4.4 节 \\
哈希 & 2. 可变时延指令 / 比较 & 不存在 & FIPS 202 & 长度驱动 & 第 4.4 节 \\
哈希 & 3. 数据相关循环 & 不存在 & FIPS 202 & 固定轮数 & 第 4.4 节 \\
哈希 & 4. 数据相关访存 & 不存在 & FIPS 202 & 常量下标 & 第 4.4 节 \\
哈希 & 5. 失败路径 & 不存在 & FIPS 202 & 不返回状态 & 第 4.4 节 \\
哈希 & 6. 异常处理 & 不存在 & FIPS 202 & 无 abort & 第 4.5 节 \\
\hline
重加密 & 1. 敏感分支 / 提前返回 & 不存在 & poly\_frommsg & 唯一跳转为循环回边 & 第 4.4 节 \\
重加密 & 2. 可变时延指令 / 比较 & 不存在 & 压缩函数 & 乘加/移位 & 第 4.4 节 \\
重加密 & 3. 数据相关循环 & 不存在 & poly\_frommsg & 32 次定长 & 第 4.4 节 \\
重加密 & 4. 数据相关访存 & 不存在 & poly\_frommsg & 固定顺序 & 第 4.4 节 \\
重加密 & 5. 失败路径 & 不存在 & indcpa\_enc & 仅分配清理 & 第 4.4 节 \\
重加密 & 6. 异常处理 & \textcolor{gray}{不适用} & — & 无异常机制 & 第 4.5 节 \\
\hline
密文比较 & 1. 敏感分支 / 提前返回 & 不存在 & ct\_memcmp & 循环体内无跳转 & 第 4.4 节 \\
密文比较 & 2. 可变时延指令 / 比较 & 不存在 & ct\_memcmp & 无分支位归约 & 第 4.4 节 \\
密文比较 & 3. 数据相关循环 & 不存在 & ct\_memcmp & 常数 1088 & 第 4.4 节 \\
密文比较 & 4. 数据相关访存 & 不存在 & ct\_memcmp & 顺序 64 字节块 & 第 4.4 节 \\
密文比较 & 5. 失败路径 & 不存在 & ct\_memcmp & 返回无构造分支 & 第 4.4 节 \\
密文比较 & 6. 异常处理 & \textcolor{gray}{不适用} & — & 不可失败 & 第 4.5 节 \\
\hline
隐式拒绝 & 1. 敏感分支 / 提前返回 & \textcolor{gray}{公开例外} & check\_sk & \texttt{call} $\to$ \texttt{jne} 校验 & 第 4.4 节 \\
隐式拒绝 & 2. 可变时延指令 / 比较 & 不存在 & ct\_cmov\_zero & 无分支掩码选择 & 第 4.4 节 \\
隐式拒绝 & 3. 数据相关循环 & 不存在 & KDF / cmov & 固定轮数与 32 字节 & 第 4.4 节 \\
隐式拒绝 & 4. 数据相关访存 & 不存在 & KDF / cmov & 固定偏移 & 第 4.4 节 \\
隐式拒绝 & 5. 失败路径 & 不存在 & decaps & 成功与失败同路径 & 第 4.4 节 \\
隐式拒绝 & 6. 异常处理 & \textcolor{gray}{不适用} & — & 无异常机制 & 第 4.5 节 \\
\hline
\multicolumn{6}{|l|}{M5 小计：存在 0 项（公开例外 1 项）、不存在 25 项、不适用 4 项、未检查 0 项（共 30 格）} \\
\hline
\end{tabularx}
}
\apptabnote{注：M5 隐式拒绝中的分支为对私钥中公开哈希校验和的合法前置检查，不接触任何私钥秘密数据。}
\end{center}

\vspace{0.8em}
\noindent\textbf{超额成果说明：}在 90 格全流程排查中，除准确锁定已知漏洞外，本作品在 M2 官方修复版本中额外排查出与 M1 同源的消息编码条件分支（重加密步骤，敏感分支类），证明了上游除法补丁未能消除该潜在分支，展现了全流程矩阵排查相比局部经验修复的系统性优势。

\fussy
